\documentclass[11pt]{article}

\usepackage[margin=1in]{geometry}
\usepackage{fontspec}
\usepackage{xeCJK}
\usepackage{xcolor}
\usepackage{hyperref}
\usepackage{graphicx}
\usepackage{array}
\usepackage{longtable}
\usepackage{listings}
\usepackage{amsmath}
\usepackage{unicode-math}

\setmainfont{Times New Roman}
\setsansfont{Arial}
\setmonofont{Consolas}
\setCJKmainfont{Microsoft YaHei}
\setCJKsansfont{Microsoft YaHei}
\setCJKmonofont{Microsoft YaHei}
\setmathfont{Cambria Math}

% 为了保证随软件分发时可以被内置 tectonic 离线渲染，本手册不强制加载
% amssymb/amsfonts。AMS 扩展符号仍在下方符号表中列出，符号预览用 Unicode
% 文字符号呈现，LaTeX 写法列给出用户在完整发行版中可使用的命令。

\hypersetup{
  colorlinks=true,
  linkcolor=blue,
  urlcolor=blue,
  citecolor=blue
}

\definecolor{GuaxiBlue}{HTML}{2563EB}
\definecolor{GuaxiSoft}{HTML}{EFF6FF}
\definecolor{GuaxiBorder}{HTML}{BFDBFE}
\definecolor{CodeBg}{HTML}{F8FAFC}

\newtheorem{definition}{Definition}
\newcommand{\vect}[1]{\mathbf{#1}}
\newcommand{\term}[1]{\textbf{\textcolor{GuaxiBlue}{#1}}}
\newcommand{\tagpill}[1]{\fbox{\textsf{\small #1}}}
% 这里使用上下边框实现提示块，避免把 \fcolorbox 的参数跨越
% \begin/\end 环境展开；这种写法对 tectonic 更稳定，也便于作为语法示例阅读。
\newenvironment{guaxicallout}[1]
  {\par\medskip\noindent\textcolor{GuaxiBorder}{\rule{\linewidth}{0.6pt}}\par
   \noindent\textbf{\textcolor{GuaxiBlue}{#1}}\par\smallskip\noindent}
  {\par\noindent\textcolor{GuaxiBorder}{\rule{\linewidth}{0.6pt}}\par\medskip}

\lstset{
  basicstyle=\ttfamily\small,
  backgroundcolor=\color{CodeBg},
  frame=single,
  breaklines=true,
  columns=fullflexible,
  keepspaces=true,
  showstringspaces=false
}

\title{\textbf{LaTeX Syntax Manual}\\LaTeX 语法查询手册}
\author{Guaxi}
\date{\today}

\begin{document}
\maketitle

\begin{abstract}
这是一份可直接被 Guaxi 内置 tectonic 引擎渲染的 LaTeX 语法手册。
It is a bilingual, renderable reference document covering mixed Chinese and
English text, page layout, fonts, table of contents, headings, paragraphs,
lists, tables, math, code, figures, links, quotes, footnotes, references,
macros, theorem-like blocks, and bibliography syntax.
\end{abstract}

\tableofcontents
\newpage

\section{文档结构 / Document Structure}
\label{sec:structure}

LaTeX documents usually contain a preamble and a document body. 中文写作时常常
需要同时处理 CJK 字体、英文术语、数学公式和代码片段；这份示例用
\texttt{fontspec} 和 \texttt{xeCJK} 设置字体，用 \texttt{geometry} 设置版面。

\begin{guaxicallout}{设计目标 / Design Goal}
本手册既是语法速查，也是渲染测试文档。The left pane shows source code, while
the right pane renders the PDF with the same preview engine used by normal
\texttt{.tex} files.
\end{guaxicallout}

\subsection{页面、字体与颜色 / Page, Font And Color}

\begin{lstlisting}[language=TeX]
\usepackage[margin=1in]{geometry}
\usepackage{fontspec}
\usepackage{xeCJK}
\setmainfont{Times New Roman}
\setCJKmainfont{Microsoft YaHei}
\definecolor{GuaxiBlue}{HTML}{2563EB}
\end{lstlisting}

\section{文本语法 / Text Markup}

Text can be \textbf{bold}, \emph{emphasized}, \texttt{monospace}, underlined by
custom packages if needed, and mixed with 中文内容。Use blank lines to start a
new paragraph.

这里是一段较长的中文段落，用于观察行距、断行、字体回退和中英文混排。
Guaxi 的 LaTeX 手册窗口会将示例源码传给内置 tectonic 引擎，并在右侧显示
PDF 预览。English technical terms like \term{compiler}, \term{package}, and
\term{macro} can appear naturally inside Chinese paragraphs.

Footnotes are written inline.\footnote{脚注支持中文和 English text. They are
rendered at the bottom of the page.}

\begin{quote}
引用块适合记录摘录、重要说明或论文中的短句。A quote environment keeps the
content visually separate from the main paragraph.
\end{quote}

\section{链接与交叉引用 / Links And Cross References}

External link example: \href{https://www.latex-project.org/}{LaTeX Project}.
Section reference example: see Section~\ref{sec:structure}. Figure reference:
Figure~\ref{fig:placeholder}. Table reference: Table~\ref{tab:syntax}.

\section{列表 / Lists}

\subsection{无序列表 / Itemize}
\begin{itemize}
  \item 中文项目：可以包含 \textbf{强调}、\texttt{代码} 和行内公式 \(a^2+b^2=c^2\)。
  \item English item: lists are useful for steps, notes, and feature summaries.
  \item Nested list:
    \begin{itemize}
      \item 子项目 / nested item
      \item \tagpill{draft} \tagpill{review} \tagpill{final}
    \end{itemize}
\end{itemize}

\subsection{有序列表与术语列表 / Enumerate And Description}
\begin{enumerate}
  \item Open a \texttt{.tex} file.
  \item Click the book icon on the toolbar.
  \item Click a source line to locate the rendered PDF preview.
\end{enumerate}

\begin{description}
  \item[Command] A control sequence such as \verb|\section|.
  \item[Environment] A block such as \verb|itemize|, \verb|tabular|, or \verb|equation|.
  \item[Package] A reusable extension loaded with \verb|\usepackage|.
\end{description}

\section{表格 / Tables}

\begin{longtable}{>{\raggedright\arraybackslash}p{0.22\linewidth}
                  >{\raggedright\arraybackslash}p{0.28\linewidth}
                  >{\raggedright\arraybackslash}p{0.38\linewidth}}
\caption{Common LaTeX syntax items / 常见 LaTeX 语法项目}
\label{tab:syntax}\\
\hline
\textbf{Feature} & \textbf{Syntax} & \textbf{Purpose} \\
\hline
\endfirsthead
\hline
\textbf{Feature} & \textbf{Syntax} & \textbf{Purpose} \\
\hline
\endhead
Heading & \verb|\section{Title}| & 生成层级标题和目录条目。 \\
Text & \verb|\textbf{}|, \verb|\emph{}| & Emphasis and inline styling. \\
Math & \verb|\[ x^2 \]| & Inline and display equations. \\
Figure & \verb|figure| & 图片、图形、标题和标签。 \\
Table & \verb|tabular|, \verb|longtable| & Structured data and comparison. \\
Code & \verb|lstlisting| & Render source code blocks. \\
\hline
\end{longtable}

\section{数学 / Mathematics}

Inline math uses \(x_i + y_i\), \(\alpha + \beta = \gamma\), and
\(\int_0^1 x^2\,dx = 1/3\).

Display math:
\[
  \sum_{i=1}^{n} i = \frac{n(n+1)}{2}.
\]

Numbered equation:
\begin{equation}
  E = mc^2
\end{equation}

Matrix and vector:
\[
  \begin{bmatrix}
    1 & 2 \\
    3 & 4
  \end{bmatrix}
  \begin{bmatrix}
    1 \\
    0
  \end{bmatrix}
  =
  \begin{bmatrix}
    1 \\
    3
  \end{bmatrix}.
\]

Piecewise function:
\[
f(x)=
\begin{cases}
0, & x < 0,\\
x^2, & x \geq 0.
\end{cases}
\]

\subsection{数学符号总览 / Mathematical Symbol Reference}

数学符号体系非常庞大，且不同宏包还会继续扩展新的符号。本节按 LaTeX 核心和
AMS 常用数学写作场景，尽量完整列出常见符号族：希腊字母、集合、逻辑、
关系、运算、箭头、微积分、几何、定界符、字体和重音。Each row shows the
rendered symbol, the LaTeX command, and its meaning.

\subsubsection{希腊字母 / Greek Letters}

\begin{longtable}{>{$}c<{$} p{0.24\linewidth} p{0.54\linewidth}}
\hline
\text{\textbf{Symbol}} & \textbf{LaTeX} & \textbf{解释 / Meaning} \\
\hline
\endfirsthead
\hline
\text{\textbf{Symbol}} & \textbf{LaTeX} & \textbf{解释 / Meaning} \\
\hline
\endhead
\alpha & \verb|\alpha| & alpha，角度、系数、显著性水平。 \\
\beta & \verb|\beta| & beta，回归系数、特殊函数参数。 \\
\gamma, \Gamma & \verb|\gamma|, \verb|\Gamma| & gamma，小写常作常数，大写常作 Gamma 函数或群。 \\
\delta, \Delta & \verb|\delta|, \verb|\Delta| & delta，微小变化；大写表示差分或判别式。 \\
\epsilon, \varepsilon & \verb|\epsilon|, \verb|\varepsilon| & epsilon，小量、误差、邻域半径。 \\
\zeta & \verb|\zeta| & zeta，Zeta 函数或阻尼比。 \\
\eta & \verb|\eta| & eta，效率、学习率或参数。 \\
\theta, \vartheta, \Theta & \verb|\theta|, \verb|\vartheta|, \verb|\Theta| & theta，角度；大写常用于渐近复杂度。 \\
\iota & \verb|\iota| & iota，嵌入映射或索引符号。 \\
\kappa & \verb|\kappa| & kappa，曲率、条件数。 \\
\lambda, \Lambda & \verb|\lambda|, \verb|\Lambda| & lambda，特征值、波长、匿名函数；大写常表谱或逻辑。 \\
\mu & \verb|\mu| & mu，均值、测度、微单位。 \\
\nu & \verb|\nu| & nu，频率、自由度。 \\
\xi, \Xi & \verb|\xi|, \verb|\Xi| & xi，随机变量、坐标或特殊函数。 \\
\pi, \varpi, \Pi & \verb|\pi|, \verb|\varpi|, \verb|\Pi| & pi，圆周率；大写表示连乘或投影。 \\
\rho, \varrho & \verb|\rho|, \verb|\varrho| & rho，密度、相关系数。 \\
\sigma, \varsigma, \Sigma & \verb|\sigma|, \verb|\varsigma|, \verb|\Sigma| & sigma，标准差；大写表示求和。 \\
\tau & \verb|\tau| & tau，时间常数、参数。 \\
\upsilon, \Upsilon & \verb|\upsilon|, \verb|\Upsilon| & upsilon，希腊字母变量。 \\
\phi, \varphi, \Phi & \verb|\phi|, \verb|\varphi|, \verb|\Phi| & phi，角度、映射、势函数。 \\
\chi & \verb|\chi| & chi，特征函数、卡方分布。 \\
\psi, \Psi & \verb|\psi|, \verb|\Psi| & psi，波函数、特殊函数。 \\
\omega, \Omega & \verb|\omega|, \verb|\Omega| & omega，角频率；大写表示欧姆、样本空间或渐近下界。 \\
\hline
\end{longtable}

\subsubsection{集合与数域 / Sets And Number Systems}

\begin{longtable}{>{$}c<{$} p{0.26\linewidth} p{0.52\linewidth}}
\hline
\text{\textbf{Symbol}} & \textbf{LaTeX} & \textbf{解释 / Meaning} \\
\hline
\endfirsthead
\hline
\text{\textbf{Symbol}} & \textbf{LaTeX} & \textbf{解释 / Meaning} \\
\hline
\endhead
ℕ & \verb|\mathbb{N}| & natural numbers，自然数集；常用 AMS 写法为 \verb|\mathbb{N}|。 \\
ℤ & \verb|\mathbb{Z}| & integers，整数集；常用 AMS 写法为 \verb|\mathbb{Z}|。 \\
ℚ & \verb|\mathbb{Q}| & rational numbers，有理数集；常用 AMS 写法为 \verb|\mathbb{Q}|。 \\
ℝ & \verb|\mathbb{R}| & real numbers，实数集；常用 AMS 写法为 \verb|\mathbb{R}|。 \\
ℂ & \verb|\mathbb{C}| & complex numbers，复数集；常用 AMS 写法为 \verb|\mathbb{C}|。 \\
\emptyset & \verb|\varnothing|, \verb|\emptyset| & empty set，空集；\verb|\varnothing| 常见于 AMS 符号包。 \\
\in, \notin & \verb|\in|, \verb|\notin| & belongs / not belongs to，属于 / 不属于。 \\
\ni & \verb|\ni| & contains as member，包含某元素。 \\
\subset, \subseteq & \verb|\subset|, \verb|\subseteq| & subset / subset or equal，真子集 / 子集。 \\
\supset, \supseteq & \verb|\supset|, \verb|\supseteq| & superset / superset or equal，超集。 \\
⊈, ⊉ & \verb|\nsubseteq|, \verb|\nsupseteq| & not subset / not superset，非子集 / 非超集，常见于 AMS 符号包。 \\
\cup, \cap & \verb|\cup|, \verb|\cap| & union / intersection，并集 / 交集。 \\
\setminus & \verb|\setminus| & set difference，集合差。 \\
\times & \verb|\times| & Cartesian product，笛卡尔积；也表示乘号。 \\
\mathcal{P}(A) & \verb|\mathcal{P}(A)| & power set，幂集。 \\
\lvert A\rvert & \verb|\lvert A\rvert| & cardinality，集合基数或绝对值。 \\
\hline
\end{longtable}

\subsubsection{逻辑与量词 / Logic And Quantifiers}

\begin{longtable}{>{$}c<{$} p{0.26\linewidth} p{0.52\linewidth}}
\hline
\text{\textbf{Symbol}} & \textbf{LaTeX} & \textbf{解释 / Meaning} \\
\hline
\endfirsthead
\hline
\text{\textbf{Symbol}} & \textbf{LaTeX} & \textbf{解释 / Meaning} \\
\hline
\endhead
\forall & \verb|\forall| & for all，对所有。 \\
\exists & \verb|\exists| & exists，存在。 \\
∄ & \verb|\nexists| & does not exist，不存在，常见于 AMS 符号包。 \\
\neg, \lnot & \verb|\neg|, \verb|\lnot| & logical not，逻辑非。 \\
\land, \wedge & \verb|\land|, \verb|\wedge| & logical and，合取。 \\
\lor, \vee & \verb|\lor|, \verb|\vee| & logical or，析取。 \\
\oplus & \verb|\oplus| & exclusive or / direct sum，异或或直和。 \\
\Rightarrow, \implies & \verb|\Rightarrow|, \verb|\implies| & implies，推出。 \\
\Leftarrow & \verb|\Leftarrow| & is implied by，被推出。 \\
\Leftrightarrow, \iff & \verb|\Leftrightarrow|, \verb|\iff| & if and only if，当且仅当。 \\
\top & \verb|\top| & true / top element，真或最大元。 \\
\bot & \verb|\bot| & false / bottom element，假或最小元。 \\
∴ & \verb|\therefore| & therefore，因此，常见于 AMS 符号包。 \\
∵ & \verb|\because| & because，因为，常见于 AMS 符号包。 \\
\hline
\end{longtable}

\subsubsection{关系符号 / Relation Symbols}

\begin{longtable}{>{$}c<{$} p{0.26\linewidth} p{0.52\linewidth}}
\hline
\text{\textbf{Symbol}} & \textbf{LaTeX} & \textbf{解释 / Meaning} \\
\hline
\endfirsthead
\hline
\text{\textbf{Symbol}} & \textbf{LaTeX} & \textbf{解释 / Meaning} \\
\hline
\endhead
= & \verb|=| & equality，等于。 \\
\ne, \neq & \verb|\ne|, \verb|\neq| & not equal，不等于。 \\
<, > & \verb|<|, \verb|>| & less than / greater than，小于 / 大于。 \\
\le, \leq & \verb|\le|, \verb|\leq| & less than or equal，小于等于。 \\
\ge, \geq & \verb|\ge|, \verb|\geq| & greater than or equal，大于等于。 \\
\ll, \gg & \verb|\ll|, \verb|\gg| & much less / much greater，远小于 / 远大于。 \\
\approx & \verb|\approx| & approximately equal，约等于。 \\
\sim & \verb|\sim| & similar / asymptotic / distributed as，相似、渐近或服从分布。 \\
\simeq & \verb|\simeq| & similar or equal，相似等价。 \\
\cong & \verb|\cong| & congruent，同余或全等。 \\
\equiv & \verb|\equiv| & equivalent / congruent，恒等或同余。 \\
\propto & \verb|\propto| & proportional to，正比于。 \\
\prec, \preceq & \verb|\prec|, \verb|\preceq| & precedes / precedes or equal，序关系。 \\
\succ, \succeq & \verb|\succ|, \verb|\succeq| & succeeds / succeeds or equal，序关系。 \\
\parallel & \verb|\parallel| & parallel，平行。 \\
\perp & \verb|\perp| & perpendicular / independent，垂直或独立。 \\
\mid, \not\mid & \verb|\mid|, \verb|\nmid| & divides / does not divide，整除 / 不整除；\verb|\nmid| 常见于 AMS 符号包。 \\
\hline
\end{longtable}

\subsubsection{运算符 / Operators}

\begin{longtable}{>{$}c<{$} p{0.26\linewidth} p{0.52\linewidth}}
\hline
\text{\textbf{Symbol}} & \textbf{LaTeX} & \textbf{解释 / Meaning} \\
\hline
\endfirsthead
\hline
\text{\textbf{Symbol}} & \textbf{LaTeX} & \textbf{解释 / Meaning} \\
\hline
\endhead
+, - & \verb|+|, \verb|-| & addition / subtraction，加 / 减。 \\
\pm, \mp & \verb|\pm|, \verb|\mp| & plus-minus / minus-plus，正负 / 负正。 \\
\cdot & \verb|\cdot| & centered multiplication dot，点乘或乘法点。 \\
\times & \verb|\times| & multiplication / Cartesian product，乘号或笛卡尔积。 \\
\div & \verb|\div| & division，除号。 \\
\ast, \star & \verb|\ast|, \verb|\star| & asterisk/star operator，星号运算。 \\
\circ & \verb|\circ| & composition / circle operator，复合运算。 \\
\bullet & \verb|\bullet| & bullet operator，点状运算符。 \\
\oplus & \verb|\oplus| & direct sum / xor，直和或异或。 \\
\ominus & \verb|\ominus| & circled minus，圈减。 \\
\otimes & \verb|\otimes| & tensor product，张量积。 \\
⊘ & \verb|\oslash| & circled slash，圈斜杠，常见于 AMS 符号包。 \\
\odot & \verb|\odot| & circled dot，圈点。 \\
\bigcirc & \verb|\bigcirc| & large circle，大圆运算符。 \\
\diamond & \verb|\diamond| & diamond operator，菱形运算符。 \\
\cap, \cup & \verb|\cap|, \verb|\cup| & intersection / union，交 / 并。 \\
\vee, \wedge & \verb|\vee|, \verb|\wedge| & logical or / and，或 / 与。 \\
\hline
\end{longtable}

\subsubsection{大型运算与微积分 / Large Operators And Calculus}

\begin{longtable}{>{$}c<{$} p{0.28\linewidth} p{0.50\linewidth}}
\hline
\text{\textbf{Symbol}} & \textbf{LaTeX} & \textbf{解释 / Meaning} \\
\hline
\endfirsthead
\hline
\text{\textbf{Symbol}} & \textbf{LaTeX} & \textbf{解释 / Meaning} \\
\hline
\endhead
\sum_{i=1}^{n} & \verb|\sum_{i=1}^{n}| & summation，求和。 \\
\prod_{i=1}^{n} & \verb|\prod_{i=1}^{n}| & product，连乘。 \\
\coprod & \verb|\coprod| & coproduct，余积。 \\
\int_a^b & \verb|\int_a^b| & definite integral，定积分。 \\
\iint_D & \verb|\iint_D| & double integral，二重积分。 \\
\iiint_V & \verb|\iiint_V| & triple integral，三重积分。 \\
\oint_C & \verb|\oint_C| & contour integral，闭合曲线积分。 \\
\lim_{x\to 0} & \verb|\lim_{x\to 0}| & limit，极限。 \\
\limsup & \verb|\limsup| & limit superior，上极限。 \\
\liminf & \verb|\liminf| & limit inferior，下极限。 \\
\frac{dy}{dx} & \verb|\frac{dy}{dx}| & derivative，导数。 \\
\partial & \verb|\partial| & partial derivative，偏导符号。 \\
\nabla & \verb|\nabla| & gradient / nabla，梯度或 nabla 算子。 \\
\Delta & \verb|\Delta| & Laplacian / difference，拉普拉斯或差分。 \\
\infty & \verb|\infty| & infinity，无穷大。 \\
\hline
\end{longtable}

\subsubsection{箭头 / Arrows}

\begin{longtable}{>{$}c<{$} p{0.28\linewidth} p{0.50\linewidth}}
\hline
\text{\textbf{Symbol}} & \textbf{LaTeX} & \textbf{解释 / Meaning} \\
\hline
\endfirsthead
\hline
\text{\textbf{Symbol}} & \textbf{LaTeX} & \textbf{解释 / Meaning} \\
\hline
\endhead
\leftarrow, \rightarrow & \verb|\leftarrow|, \verb|\rightarrow| & left / right arrow，左箭头 / 右箭头。 \\
\leftrightarrow & \verb|\leftrightarrow| & bidirectional arrow，双向箭头。 \\
\Leftarrow, \Rightarrow & \verb|\Leftarrow|, \verb|\Rightarrow| & double arrows，双线箭头。 \\
\Leftrightarrow & \verb|\Leftrightarrow| & double bidirectional arrow，双线双向箭头。 \\
\mapsto & \verb|\mapsto| & maps to，映射到。 \\
\longmapsto & \verb|\longmapsto| & long maps to，长映射箭头。 \\
\uparrow, \downarrow & \verb|\uparrow|, \verb|\downarrow| & up / down arrows，上 / 下箭头。 \\
\updownarrow & \verb|\updownarrow| & up-down arrow，上下双向箭头。 \\
\nearrow, \searrow & \verb|\nearrow|, \verb|\searrow| & diagonal arrows，东北 / 东南箭头。 \\
\swarrow, \nwarrow & \verb|\swarrow|, \verb|\nwarrow| & diagonal arrows，西南 / 西北箭头。 \\
\hookrightarrow & \verb|\hookrightarrow| & injection arrow，嵌入或单射。 \\
↠ & \verb|\twoheadrightarrow| & surjection arrow，满射，常见于 AMS 符号包。 \\
\hline
\end{longtable}

\subsubsection{定界符与括号 / Delimiters}

\begin{longtable}{>{$}c<{$} p{0.28\linewidth} p{0.50\linewidth}}
\hline
\text{\textbf{Symbol}} & \textbf{LaTeX} & \textbf{解释 / Meaning} \\
\hline
\endfirsthead
\hline
\text{\textbf{Symbol}} & \textbf{LaTeX} & \textbf{解释 / Meaning} \\
\hline
\endhead
(x) & \verb|(x)| & parentheses，圆括号。 \\
\lbrack x\rbrack & \verb|\lbrack x\rbrack| & square brackets，方括号。 \\
\{x\} & \verb|\{x\}| & braces，花括号。 \\
\langle x\rangle & \verb|\langle x\rangle| & angle brackets，尖括号。 \\
\lvert x\rvert & \verb|\lvert x\rvert| & absolute value，绝对值。 \\
\lVert x\rVert & \verb|\lVert x\rVert| & norm，范数。 \\
\lfloor x\rfloor & \verb|\lfloor x\rfloor| & floor，向下取整。 \\
\lceil x\rceil & \verb|\lceil x\rceil| & ceiling，向上取整。 \\
\left(\frac{a}{b}\right) & \verb|\left( ... \right)| & auto-sized delimiters，自动伸缩定界符。 \\
\hline
\end{longtable}

\subsubsection{几何与常用特殊符号 / Geometry And Miscellaneous Symbols}

\begin{longtable}{>{$}c<{$} p{0.28\linewidth} p{0.50\linewidth}}
\hline
\text{\textbf{Symbol}} & \textbf{LaTeX} & \textbf{解释 / Meaning} \\
\hline
\endfirsthead
\hline
\text{\textbf{Symbol}} & \textbf{LaTeX} & \textbf{解释 / Meaning} \\
\hline
\endhead
\angle & \verb|\angle| & angle，角。 \\
\triangle & \verb|\triangle| & triangle，三角形。 \\
□ & \verb|\square| & square，正方形或证明结束符，常见于 AMS 符号包。 \\
◇ & \verb|\Diamond| & modal diamond / diamond，模态菱形或菱形，常见于 AMS 符号包。 \\
\forall x\in A & \verb|\forall x\in A| & quantified membership，带量词的集合命题。 \\
\ell & \verb|\ell| & script-like l，常表示直线。 \\
\Re, \Im & \verb|\Re|, \verb|\Im| & real / imaginary part，实部 / 虚部。 \\
\aleph & \verb|\aleph| & aleph number，阿列夫数。 \\
\hbar & \verb|\hbar| & reduced Planck constant，约化普朗克常数。 \\
\prime & \verb|\prime| & prime mark，撇号。 \\
\dagger, \ddagger & \verb|\dagger|, \verb|\ddagger| & dagger symbols，剑号脚注或伴随记号。 \\
\diamondsuit, \heartsuit & \verb|\diamondsuit|, \verb|\heartsuit| & suit symbols，花色符号。 \\
\clubsuit, \spadesuit & \verb|\clubsuit|, \verb|\spadesuit| & suit symbols，花色符号。 \\
\hline
\end{longtable}

\subsubsection{数学重音、字体与函数 / Accents, Fonts And Functions}

\begin{longtable}{>{$}c<{$} p{0.30\linewidth} p{0.48\linewidth}}
\hline
\text{\textbf{Symbol}} & \textbf{LaTeX} & \textbf{解释 / Meaning} \\
\hline
\endfirsthead
\hline
\text{\textbf{Symbol}} & \textbf{LaTeX} & \textbf{解释 / Meaning} \\
\hline
\endhead
\hat{x} & \verb|\hat{x}| & estimator / hat，估计量或帽号。 \\
\bar{x} & \verb|\bar{x}| & average / conjugate，均值或共轭。 \\
\tilde{x} & \verb|\tilde{x}| & approximation / equivalence，近似或等价类。 \\
\vec{x} & \verb|\vec{x}| & vector arrow，向量。 \\
\dot{x}, \ddot{x} & \verb|\dot{x}|, \verb|\ddot{x}| & first / second time derivative，一阶 / 二阶时间导数。 \\
\overline{AB} & \verb|\overline{AB}| & overline，上横线。 \\
\underline{AB} & \verb|\underline{AB}| & underline，下横线。 \\
\widehat{ABC} & \verb|\widehat{ABC}| & wide hat，宽帽号。 \\
\widetilde{ABC} & \verb|\widetilde{ABC}| & wide tilde，宽波浪号。 \\
\mathbf{x} & \verb|\mathbf{x}| & bold math，粗体数学字母。 \\
\mathit{x} & \verb|\mathit{x}| & italic math，斜体数学字母。 \\
\mathrm{d}x & \verb|\mathrm{d}x| & upright roman text，正体微分符号。 \\
\mathcal{F} & \verb|\mathcal{F}| & calligraphic letters，花体字母。 \\
ℝ & \verb|\mathbb{R}| & blackboard bold，黑板粗体；完整发行版常配合 AMS 字体使用。 \\
\sin x, \cos x & \verb|\sin x|, \verb|\cos x| & trigonometric functions，三角函数。 \\
\tan x, \cot x & \verb|\tan x|, \verb|\cot x| & tangent / cotangent，正切 / 余切。 \\
\log x, \ln x & \verb|\log x|, \verb|\ln x| & logarithms，对数。 \\
\exp x & \verb|\exp x| & exponential function，指数函数。 \\
\max, \min & \verb|\max|, \verb|\min| & maximum / minimum，最大 / 最小。 \\
\sup, \inf & \verb|\sup|, \verb|\inf| & supremum / infimum，上确界 / 下确界。 \\
\arg, \det & \verb|\arg|, \verb|\det| & argument / determinant，辐角 / 行列式。 \\
\dim, \ker & \verb|\dim|, \verb|\ker| & dimension / kernel，维数 / 核。 \\
\hline
\end{longtable}

\section{代码块 / Code Blocks}

\begin{lstlisting}[language=TeX,caption={LaTeX macro example}]
\newcommand{\vect}[1]{\mathbf{#1}}
\section{A Section}
\end{lstlisting}

\begin{lstlisting}[language=Python,caption={Python example}]
def area(radius):
    return 3.14159 * radius ** 2
\end{lstlisting}

\section{图片与图形 / Images And Figures}
\label{sec:figures}

为了保证手册离线可渲染，这里使用 \verb|\fbox| 和 \verb|\rule| 生成图片占位。
真实图片可使用 \verb|\includegraphics[width=0.7\linewidth]{path/to/image.png}|；
相对路径通常以当前 \texttt{.tex} 文件所在目录为基准。

\begin{figure}[h]
\centering
\fcolorbox{GuaxiBlue}{GuaxiSoft}{%
  \begin{minipage}{0.72\linewidth}
    \centering
    \vspace{1.1cm}
    \textbf{Generated image placeholder}\\
    这是无需外部文件的图片/图形示例
    \vspace{1.1cm}
  \end{minipage}
}
\caption{A generated image-like figure / 生成式图片占位}
\label{fig:placeholder}
\end{figure}

\section{宏、环境与定理 / Macros, Environments And Theorems}

\begin{definition}
A macro is a reusable command. This document defines \verb|\vect| and uses it as
\(\vect{x}\), \(\vect{y}\), and \(\vect{z}\).
\end{definition}

\begin{guaxicallout}{自定义环境 / Custom Environment}
\texttt{guaxicallout} is a custom environment built with a minipage and colored
frame. 自定义环境适合封装重复出现的版式块。
\end{guaxicallout}

\section{多栏、并排块与分页 / Side-by-side Blocks And Page Breaks}

\noindent
\begin{minipage}{0.47\linewidth}
\textbf{左侧 / Left}\\
Minipages are useful for side-by-side blocks, figures, forms, and summaries.
中文说明也可以放在同一个并排块中。
\end{minipage}
\hfill
\begin{minipage}{0.47\linewidth}
\textbf{右侧 / Right}\\
They keep content in a local box. Use \verb|\hfill| to separate two minipages.
\end{minipage}

\newpage

\section{附录 / Appendix}

This final section makes the document long enough for scrolling, searching, and
source-to-preview navigation. 最后一节用于测试窗口缩放、PDF 搜索、源码行点击定位等功能。

\begin{lstlisting}[language=TeX,caption={External image syntax}]
\begin{figure}
  \centering
  \includegraphics[width=0.7\linewidth]{assets/example.png}
  \caption{An external image.}
\end{figure}
\end{lstlisting}

\section{参考文献 / Bibliography}

Manual bibliography syntax can be written without an external \texttt{.bib} file.
Example citation: \cite{knuth1984texbook}.

\begin{thebibliography}{9}
\bibitem{knuth1984texbook}
Donald E. Knuth.
\emph{The TeXbook}.
Addison-Wesley, 1984.
\end{thebibliography}

\end{document}
